\section{Higher-component pencils and gluing}
\label{sec:gluing}

We retain the source and target chamber spaces $B_S,B_T$, the type
permutation $p$, and the sharp angular measurements constructed in
Section~\ref{sec:transfer}.  For a panel link, $P_z$ and $Q_z$ denote
the start and end missing-vertex measurements in their separate compact
frames, and $g$ is the target projectivity from the end pencil to the
start pencil.  The two inputs of $J$ in Equation~\eqref{eq:input-6}
remain independent tensor factors.  A cut on a vertex coordinate of a
chamber measurement means its pullback under the corresponding flag
projection.

The isolated positions, including commutation between their two normal
copies, were treated in Section~\ref{sec:locality}.  At a finite place,
Section~\ref{sec:finite} identifies each one-copy panel law with
evaluations of homogeneous analytic homeomorphisms.  We first establish
the corresponding conclusion for the non-isolated archimedean
positions.  The Taylor argument below is used only at these positions.
Afterwards the arguments joining the two copies and joining different
panels apply at all places.

\subsection{Real Taylor density for the angular pairing}

\begin{lemma}[Density of pairing powers]
\label{lem:gluing-taylor-density}
Let $\mathcal V$ be an archimedean target vertex variety and
$\mathcal V^*$ its opposite-type dual variety.  Use the nilradical
determinant-line representation and normalized absolute pairing
$k_T$ of Section~\ref{sec:heights}.  Fix $v_0\in\mathcal V^*$ and a
compact set $K_*$ of vertices transverse to $v_0$.  If $O$ is any
sufficiently small open neighborhood of $v_0$ transverse to $K_*$,
there is a countable set $\mathcal E\subset\mathbb R$ such that
\[
 \overline{\operatorname{span}}^{\|\cdot\|_\infty}
 \{u\mapsto k_T(v,u)^{\mathrm i\theta}:v\in O\}=C(K_*)
 \qquad(\theta\notin\mathcal E).
\]
The span is complex linear.
\end{lemma}

\begin{proof}
Use real coordinates $x$ on the opposite unipotent cell determined by
$v_0$, and coordinates $y$ on a dual unipotent cell with center $y=0$
at $v_0$.  Normalize the algebraic highest-line lifts against the two
fixed opposite centers.  The pairing of these lifts is a real
polynomial $F(y,x)$ with $F(0,x)=1$.  The homogeneous norms used to
define $k_T$ therefore give
\begin{equation}
 \label{eq:gluing-normalized-pairing}
 k_T(v(y),u(x))=a(y)b(x)|F(y,x)|,
 \qquad a(y)>0,\quad b(x)>0.
\end{equation}
After shrinking $O$, the polynomial $F$ is positive on a neighborhood
of $\{0\}\times K_*$.  Thus its complex powers there use the branch
defined by the real logarithm.

A central parabolic cocharacter gives positive integral weights to
the nilpotent coordinates.  Multiplying the cocharacter if necessary
makes all weights integral.  Write $\delta_t$ for the corresponding
graded dilation.  The cocharacter acts oppositely on the two cells, so
\[
 F(\delta_t y,\delta_{t^{-1}}x)=F(y,x).
\]
For a multi-index $\alpha$ put
\[
 c_\alpha(s,x)=\frac{1}{\alpha!}
       \partial_y^\alpha F(y,x)^s\big|_{y=0},\qquad s\in\mathbb C.
\]
If $|\alpha|_{\mathrm{gr}}=d$, then $c_\alpha(s,x)$ is a weighted
homogeneous polynomial of degree $d$ in $x$.  Its coefficients are
polynomial in $s$: in the finite Taylor jet at $y=0$, the expansion
of $(1+(F-1))^s$ uses only finitely many falling factorials in $s$.
In particular, for each $d$ only finitely many such coefficients
occur.

We prove that these coefficients span all polynomials of degree $d$
for a generic $s$.  Complexify the real determinant-line cyclic
representation.  The wedge of the complexified parabolic nilradical
is a highest vector, of weight the sum of its positive roots.
Complete reducibility makes the module generated by that vector
one irreducible highest-weight module; it is this cyclic module
that we complexify.  These are the usual highest-weight and
parabolic constructions \cite{Humphreys1972,Borel1991}.
As in Section~\ref{sec:links}, its highest-line
orbit is the complex parabolic variety $X=G_{\mathbb C}/P_{\mathbb C}$,
and the highest weight is strictly parabolic: its zero simple-coroot
pairings are exactly those indexed by the Levi of $P_{\mathbb C}$.
This assertion also covers a complex group regarded as real; its
complexification and parabolic can have several factors.

Let $h$ be the highest coordinate paired with the dual center.  The
principal open subset $\{h\ne0\}$ of $X$ is precisely the opposite
unipotent cell.  Indeed, on an opposite Bruhat cell indexed by a Weyl
element $w$, the highest coordinate can be nonzero only when $w$
fixes the highest weight.  Strict parabolicity makes this equivalent
to $w\in W_{P_{\mathbb C}}$, which is the open-cell index.  Conversely
the normalized unipotent lift has $h=1$ throughout that cell.
The cell is algebraically isomorphic to its unipotent group, hence
to affine space in the coordinates $x$.

The coordinate ring on this principal open is the degree-zero part
of the homogeneous coordinate ring localized at $h$.  Consequently
every polynomial in $x$ is the restriction of $q/h^n$ for some
homogeneous coordinate polynomial $q$ of degree $n$.  Given a basis
of the finite-dimensional space $\mathcal R_d$ of weighted degree
$d$ polynomials, choose one common $n$ for all these expressions by
multiplying numerators and denominators by powers of $h$.
The same expressions are then available at every larger $n$.

For integral $n$, the line tensors of order $n$ span the irreducible
Cartan component of highest weight $n\lambda$.  The dual highest-line
tensors span its dual.  They still span when their parameters are
restricted to an arbitrary real open dual cell: a linear functional
vanishing there is a polynomial vanishing on a real open subset of
the real coordinate space, and thus on its complexification.
After dehomogenization, it follows that the functions
$F(y,\cdot)^n$, with $y$ in that real open cell, span all restrictions
of homogeneous coordinate polynomials of degree $n$.  Taking the
weighted degree $d$ part gives
\[
 \mathcal R_d
   =\operatorname{span}\{c_\alpha(n,\cdot):
                                      |\alpha|_{\mathrm{gr}}=d\}
 \quad\text{for all sufficiently large integers }n.
\]
Here the equality follows because matching graded degrees makes the
degree $d$ part of $F(y,x)^n$ exactly
$\sum_{|\alpha|_{\mathrm{gr}}=d}y^\alpha c_\alpha(n,x)$.

In a fixed monomial basis of $\mathcal R_d$, choose a maximal minor
of the coefficient matrix of the $c_\alpha(s,\cdot)$ which is
nonzero at one such integer $n$.  It is a nonzero polynomial in
$s$, so it vanishes at only finitely many complex parameters.
The union of these finite exceptional sets over all $d$ is
countable.  For $s=\mathrm i\theta$ outside this union, every
polynomial in $x$ is a linear combination of Taylor coefficients
$c_\alpha(s,\cdot)$.

Taylor differentiation at $y=0$ is the uniform limit on $K_*$ of
ordinary real difference quotients, because $F$ stays positive on
a neighborhood of $\{0\}\times K_*$.  Each coefficient therefore
belongs to the uniformly closed span of $F(y,\cdot)^s$ for
arbitrarily small real $y$.  Real-coordinate polynomials, with
complex coefficients, contain constants, separate points of $K_*$,
and are closed under conjugation.  The Stone--Weierstrass Theorem
makes their uniform closure $C(K_*)$.  Finally $a(y)^s$ in
Equation~\eqref{eq:gluing-normalized-pairing} is a nonzero scalar
for each $y$, and multiplication by the fixed continuous
nonvanishing function $b(x)^s$ is an automorphism of $C(K_*)$.
These norm factors do not change density.  Intersecting the
exceptional complex parameters with the imaginary axis gives
$\mathcal E$.
\end{proof}

\Needspace{6\baselineskip}
\subsection{Reduction of a real panel link}

\begin{lemma}[Separated projectivity blocks]
\label{lem:gluing-separated-blocks}
Consider a one-leg link omitting a single non-isolated
archimedean type.  Fix compact endpoint frames and two sufficiently
small compact output windows $U,L$ about distinct projectivities
with the same endpoint panels.  Let $E_U,E_L$ be their output
spectral projections.  There is an ordinary dual-vertex
neighborhood $O$ such that, for generic source chamber labels $Z$,
\begin{equation}
 \label{eq:input-26}
 E_UJ\bigl(1\otimes P_Z(O)X P_Z(O)\bigr)J^*E_L=0
 \qquad\text{for every bounded }X.
\end{equation}
For each fixed completion $z$, one can shrink $O$ and find a
compact set $K_*$ of possible start values, transverse to all of
$O$, such that
\begin{equation}
 \label{eq:gluing-phase-block}
 E_UJ\bigl(1\otimes
     k_T(v,P_z)^{\mathrm i\theta}b_{F,i}^{\mathrm i\theta}
                 \bigr)J^*E_L=0
 \qquad(v\in O,\ \theta\in\mathbb R).
\end{equation}
Here $F$ is a tuple in the completed chamber containing the missing
vertex.  The block in Equation~\eqref{eq:gluing-phase-block} is
unchanged by restricting the $P_z$ value to $K_*$, and its
complement vanishes also with additional own-$P_z$ multipliers.
\end{lemma}

\begin{proof}
Two distinct projectivities have different inverse images at some
point of the start pencil.  Choose an ordinary dual vertex $v_0$
whose transverse projection onto that pencil is such a point.
Shrink a neighborhood $O$ of $v_0$ and the windows $U,L$, including
their endpoint base windows, so that the two inverse-image sets
remain separated with positive slack.  Distinct base values
already give zero blocks by sharpness and their commutation with
the operators under consideration, so only a common base chart
needs treatment.

For a generic source chamber $Z$, project its pertinent dual vertex
onto the source start panel and denote the resulting completion
by $z_Z$.  On the ordered start cut $P_Z(O)$ the value at
$P_{z_Z}$ is the corresponding transverse projection.  To see the
operator statement, insert a gallery shortened from the longest
word.  The angular nonopposition rule of Section~\ref{sec:transfer}
forces its endpoint to be the unique completion of the panel
nonopposite to that ordinary transverse dual vertex.  The
multiplied gallery identity gives this substitution together with
the sharp base cuts.

Apply Equation~\eqref{eq:input-6} after that substitution.  On the
left output window $U$, the independent end value $Q_{z_Z}$ is
forced into the corresponding inverse-image cut; on the right
window $L$, the adjoint argument forces it into the disjoint
inverse-image cut.  For precision, partition the compact base
windows and the projected completion windows into sufficiently
small patches, insert their sharp base and $P_{z_Z}$ projections
next to $J$ on the appropriate side of $P_Z(O)$, and insert the
independent $Q_{z_Z}$ projection there.  All these moves involve
either a completion and its own base or separate tensor factors.
The two final end cuts are orthogonal and act in the tensor factor
untouched by $X$.  Their product is zero for every $X$; summing
the finite partitions proves Equation~\eqref{eq:input-26}.
Compact exhaustion removes the preliminary frame restrictions.

Now compress Equation~\eqref{eq:input-23} with the isometry $V$
of Equation~\eqref{eq:input-22}, testing ordinary chamber vectors
whose scalar supports are in $O$.  On the source side the
integrands are flanked by $P_Z(O)$: this follows directly from
the formula $V(f\otimes\xi)(Z)=
c_0^{-1/2}D_Z^{-1/2}P_Z(f)\xi$ and the own-coordinate
commutations established in Section~\ref{sec:heights}.
Initially cut $D_Z^{-1/2}$ by bounded spectral truncations.
Equation~\eqref{eq:input-26} then annihilates every integrand.
The truncated images converge in their direct-integral Hilbert
norm to the images under $V$.  Pairing two such images with the
bounded multiplier in Equation~\eqref{eq:input-23} passes to the
limit by Cauchy--Schwarz.  Thus the compressed operators converge
ultraweakly, which suffices after multiplication by the bounded
operators $E_UJ$ and $J^*E_L$.  Arbitrary scalar tests and Hilbert
vectors supported in $O$ now give
Equation~\eqref{eq:gluing-phase-block} for almost every ordinary
dual vertex $v\in O$.

The multiplier between $J,J^*$ commutes with the sharp starting
panel, its own $P_z$, and the independent $Q_z$.  Hence only common
data for these slots can contribute to the block.  By
Equation~\eqref{eq:input-6}, a possible start value must have a
common inverse image under projectivities in $U$ and $L$.
The separated-inverse-image choice excludes a neighborhood of the
completion projected from $v_0$.  Within the start panel this is
exactly the locus which can fail to be transverse to $v_0$.
After shrinking the base and projectivity windows, all possible
common start values therefore lie in a compact set $K_*$
transverse to $v_0$.  Compact partitions in the commuting own
slots show that every complementary contribution vanishes,
also after inserting an additional own-$P_z$ multiplier.
This uses no arbitrary restriction inside $K_*$.

Shrink $O$ so that $K_*\times O$ is transverse.  On this compact
set the function $v\mapsto k_T(v,\cdot)^{\mathrm i\theta}$ is
continuous in the uniform norm.  The almost-everywhere zero
blocks consequently extend to every $v$ in the smaller $O$,
proving the claim.
\end{proof}

\begin{proposition}[One-copy commutation in a real pencil]
\label{prop:gluing-one-pencil}
The output projectivity PVM reduces $\ran J$.  On the input
space its reduced calculus commutes with every $P_z$ and $Q_z$,
and Equation~\eqref{eq:input-6} becomes the joint spectral identity
$P_z=gQ_z$.  All $P_z$ in the pencil commute mutually, and all
$Q_z$ in the pencil commute mutually.
\end{proposition}

\begin{proof}
Fix the separated windows of
Lemma~\ref{lem:gluing-separated-blocks}.  For
$\theta\notin\mathcal E$, Lemma~\ref{lem:gluing-taylor-density}
and the compact restriction to $K_*$ replace the pairing power
in Equation~\eqref{eq:gluing-phase-block} by any continuous
function $f(P_z)$.  Uniform approximation is legitimate because
$b_{F,i}^{\mathrm i\theta}$ is unitary and $J$ is an isometry.
Choose nonexceptional $\theta_n\to0$.  Strong continuity of this
unitary group, with all angular factors bounded, gives
\[
 E_UJJ^*E_L=0,
 \qquad
 E_UJ(1\otimes P_z(f))J^*E_L=0.
\]
The compact restriction in the preceding lemma has removed all
other contributions, so these equations hold for the full blocks.
The argument with the start and end slots interchanged gives the
same assertion for end multipliers.

Every two distinct projectivity values, and every two distinct
base values, have separating windows of this kind.  Exhaust by
countably many compact charts.  The vanishing of all these
off-diagonal blocks says precisely that $JJ^*$ and the displayed
compressed operators commute with the output PVM.  One can verify
this criterion first for disjoint compact spectral sets, covering
their product by finitely many separating rectangles, and then
use spectral regularity and increasing exhaustion for arbitrary
Borel cuts.  Thus $\ran J$ reduces the output PVM.

Pull its restriction back by $J$.  All own input angular operators
commute with this reduced projectivity calculus.  Consequently
the graph substitution of Equation~\eqref{eq:input-6}, first on
compact windows and then by exhaustion, reads $P_z=gQ_z$ in
joint calculus.  For any $w$, the operator $P_w(f)$ commutes with
both $g$ and $Q_z$, the latter because it acts on the other input
factor.  It therefore commutes with $P_z(f')$.  This proves
mutual commutation of the start family.  Interchanging the two
inputs proves the end assertion.
\end{proof}

\subsection{The homogeneous form of a real pencil map}

\begin{proposition}[Projectivity-orbit classification]
\label{prop:gluing-homogeneous-pencil}
On almost every sharp endpoint fibre, the one-copy real pencil
measurements are evaluations of random homogeneous
diffeomorphisms from the source pencil onto the target pencil.
The maps and their inverses have jointly measurable versions as
the endpoint fibres vary.  Their evaluations recover the
original PVMs at every fixed source label.
\end{proposition}

\begin{proof}
Fix a good pair of sharp endpoint fibres in
Equation~\eqref{eq:input-6}.  The disintegration in
Section~\ref{sec:transfer}, with a genuine continuous source lift
when necessary, supplies strongly continuous source Levi
representations on these fibres, transporting the labels while
fixing the bases.  Write $\mathcal B_S,\mathcal B_T$ for the
source and target pencil boundaries, and $S,T$ for their
effective panel action groups.  The identity action groups have
noncompact simple real rank-one Lie algebras
$\mathfrak s,\mathfrak t$.  Compact normal factors and split
centers are ineffective; finite components of the target
projectivity group can be retained.

Choose faithful normal states on the separable abelian angular
algebras of the two fibres.  Their scalar spectral laws realize
the evaluations as random measurable maps
$\phi,\psi:\mathcal B_S\to\mathcal B_T$, initially modulo
ordinary source null sets.  Here is a useful measurability
construction.  Strong source continuity gives, for a countable
family of continuous target functions, continuity of their
evaluations into scalar convergence in measure.  Finite
measurable approximations of the source label, chosen with
summable errors in measure on compact pieces, give jointly
measurable limits.  The countably many target relations identify
a target point almost everywhere.  Thus the map classes are
random variables in the Polish space
$L^0(\mathcal B_S,\mathcal B_T)$.

Use the product of the two faithful states.  The start and end
map variables are then independent.  The reduced projectivity
calculus of Proposition~\ref{prop:gluing-one-pencil} witnesses,
by Fubini, the relation $\phi=g\psi$ as map classes for
product-almost every pair.  The projectivity-orbit relation is
Borel.  Indeed, for each compact subset $K\subset T$, the set
\[
 \{(\phi,\psi):\phi=t\psi\text{ for some }t\in K\}
\]
is closed: take a convergent subsequence of the compact parameters
and use continuity of postcomposition in measure.  A countable
compact exhaustion of $T$ gives the claim.  Fubini now supplies
a representative $\phi_0$ such that both independent map laws
are concentrated on the one orbit $T\phi_0$.

This representative is not essentially constant.  Otherwise all
maps in the orbit would be constant, and source continuity would
make every $P_z$ the same PVM on the fibre.  Choose a nonconstant
polynomial coordinate test $f$ on the ordinary incident-line
embedding, a nonzero vector $\xi$, and put
$\eta=P_{z_0}(f)\xi$.  We would have
$P_z(f)\xi-\eta=0$ for all $z$.  Fibre polynomial fullness from
Section~\ref{sec:transfer}, after positive homogeneous
normalization, would give $f(u)\xi-\eta=0$ for every ordinary
target residue value $u$.  Two values of $f$ contradict
$\xi\ne0$.  Equivalently, a constant measured line cannot
contain the tensor product of the fibre with the full
ordinary incident-line span, which has dimension greater than
one.

For each $s\in S$, source covariance is implemented unitarily on
the fibre.  It preserves the faithful spectral measure class,
whether or not it preserves the chosen state.  Thus the law
after precomposition by $s$ is still concentrated on
$T\phi_0$.  Precomposition commutes with the target action;
intersecting the two full-measure orbit sets shows that
$\phi_0\circ s\in T\phi_0$.  Consider
\begin{equation}
 \label{eq:gluing-pair-subgroup}
 L=\{(s,t)\in S\times T:
               t\phi_0=\phi_0\circ s
                  \text{ modulo ordinary null sets}\}.
\end{equation}
It is a closed subgroup projecting onto $S$.  For closedness,
target composition is continuous by compactness of
$\mathcal B_T$, and source composition is continuous in measure
by approximation with continuous functions and locally uniform
bounds on the smooth change-of-variables densities.

The Closed Subgroup Theorem makes $L$ a Lie group.  Its
projection onto $S$ is onto also on Lie algebras: otherwise
the Constant Rank Theorem on countably many Lie charts would
make its image a Haar-null subset of $S$.  The radical of
$\mathfrak l$ has zero image in the simple algebra
$\mathfrak s$.  A Levi decomposition \cite[Appendix~B, Theorem~B.2]{Knapp2023}
therefore has a simple summand mapping isomorphically onto $\mathfrak s$; distinct
simple summands have commuting ideal images, so at most one
can have a nonzero image.  In
$\mathfrak s\oplus\mathfrak t$, this summand is the graph
\[
 X\longmapsto (X,\alpha(X))
\]
of a Lie homomorphism $\alpha:\mathfrak s\to\mathfrak t$.
If $\alpha=0$, Equation~\eqref{eq:gluing-pair-subgroup} makes
$\phi_0$ invariant modulo null sets under the connected
transitive source group, hence essentially constant.  Thus
$\alpha\ne0$, and simplicity gives injectivity.  For a panel
of type $i$ this proves
$\dim\mathfrak s_i\le\dim\mathfrak s_{p(i)}$.
The archimedean types form a union of cycles of $p$ by
Section~\ref{sec:finite}; going around each finite cycle forces
equality at every step.  Hence $\alpha$ is an isomorphism.

Let $\widetilde S$ be the simply connected source cover and
$a:\widetilde S\to T^\circ$ the homomorphism integrating
$\alpha$.  It is onto because its differential is onto.
Let $P_S$ be the full preimage of the source minimal
parabolic, so $\mathcal B_S=\widetilde S/P_S$.  Equation
\eqref{eq:gluing-pair-subgroup} gives
$\phi_0(sx)=a(s)\phi_0(x)$ for every $s$ and almost every $x$.
On $\widetilde S$, the measurable function
\[
 r\longmapsto a(r)^{-1}\phi_0(rP_S)
\]
is therefore invariant modulo Haar-null sets under each left
translation.  Fubini, or convolution of its bounded scalar
coordinate tests with a compactly supported approximate identity,
makes it essentially constant, say equal to $u_0$.
Its right $P_S$ transformation shows that $a(P_S)$ fixes $u_0$.
The recovered map is consequently
\[
 rP_S\longmapsto a(r)u_0.
\]
An isomorphism of real semisimple Lie algebras carries minimal
parabolic Lie algebras to minimal parabolic Lie algebras
\cite{Borel1991}.
Thus the inclusion of $\alpha(\operatorname{Lie}P_S)$ in the
stabilizer algebra of $u_0$ is equality.  Boundary stabilizers
are the full normalizers of these parabolic algebras.  If
$a(r)$ fixes $u_0$, then $r$ normalizes
$\operatorname{Lie}P_S$, and hence $r\in P_S$.
This proves
$a^{-1}(\operatorname{Stab}_{T^\circ}(u_0))=P_S$.
The displayed map is a homogeneous diffeomorphism, rather
than a map between nontrivial oriented covers.  Central kernels
and ineffective compact factors lie in the stabilizers;
finite target components also act by diffeomorphisms.

Every almost-sure map class now has a unique continuous
homeomorphism version.  These versions are measurable without
a measurable choice of $\phi_0$ on different fibres.
The space of homeomorphisms between compact metrizable pencils,
with uniform convergence of maps and inverses, is Polish.
Its injection into $L^0(\mathcal B_S,\mathcal B_T)$ is Borel
and injective, since ordinary source measure has full support.
The Lusin--Souslin Theorem \cite[Corollary~15.2]{Kechris1995} gives a Borel inverse on its image.
Apply that inverse in countably many compact endpoint charts.
Finally, the original evaluations and these continuous
evaluations agree for almost every label.  Both are continuous
in the label in scalar convergence in measure, the former by
source transport and the latter by bounded convergence.
They therefore agree for every fixed label on all continuous
target tests, and hence as PVMs.
\end{proof}

\subsection{Commutation of the two normal copies}

\begin{lemma}[Two-copy pencil commutation]
\label{lem:gluing-two-copy}
At every non-isolated position, the two normal copies of the
one-copy pencil calculus commute with one another.  In
particular, adjacent chamber measurements commute, also when
each chamber carries both copies.
\end{lemma}

\begin{proof}
Disintegrate the common two-leg link over its joint sharp endpoint
faces. These faces commute with their respective completion observations.
Proposition~\ref{prop:gluing-homogeneous-pencil} and the finite-place
classification in Section~\ref{sec:finite} therefore give, in each copy
$j=1,2$, a start homeomorphism $B_j$ and an end homeomorphism $A_j$
from the source pencil to the corresponding target pencil. Their
continuous versions and inverses persist under normal copying and
disintegration.

For each $j$ separately, $B_j$ and $A_j$ act in independent input
factors, so their joint spectral calculus defines the target-to-target
map
\[
 C_j=B_jA_j^{-1}.
\]
Write $E_j$ for the output projectivity PVM in that coordinate of the
two-leg link. For this fixed $j$, all evaluations of $A_j$ and $B_j$
commute.
We may therefore multiply the residue graph identities in
Equation~\eqref{eq:input-6} by further cuts from this same input
calculus, imposing finitely many determining-label constraints together.
Choose the labels in a countable dense set. Continuity makes their
simultaneous graph equations $g_jA_j(z)=B_j(z)$ determine
$g_j=B_jA_j^{-1}$. Separation by compact projectivity windows, followed
by spectral exhaustion, consequently gives
\[
 E_j(H)J=J\,1_H(C_j)
\]
for every Borel projectivity cut $H$. This identity also puts the law
of $C_j$ on projectivities. The two output coordinate PVMs commute in
the common link. Their intertwining identities therefore imply
commutation of the two input composite calculi. Source transport in
the start factor gives the same conclusion for
\[
 B_1sA_1^{-1}\quad\text{and}\quad B_2sA_2^{-1},
 \qquad s\in S.
\]
This argument uses the joint calculus of $(B_j,A_j)$ separately for
each $j$; commutation of $B_1$ with $B_2$ is still to be proved.

Let $K_S$ be a compact subgroup transitive on the source
pencil boundary; at a finite place use a special maximal
compact.  Its normalized Haar measure induces the ordinary
pencil probability $m_S$.  Fix disjoint compact sets
$\mathcal U,\mathcal L$ in the start $B_1$ homeomorphism
space, with spectral projections $R_{\mathcal U},R_{\mathcal L}$.
Fix a source transport $s_n$ and a support point $a_0$ of
the end $A_1$ law.  For each common $k\in K_S$, the two
compact systems
\[
 \{b s_nk a_0^{-1}:b\in\mathcal U\},\qquad
 \{b s_nk a_0^{-1}:b\in\mathcal L\}
\]
are disjoint, since right composition is injective.
Their separation in the topology of maps and inverses is
uniform for this common $k$, by compactness of
$\mathcal U\times\mathcal L\times K_S$.
Hence a sufficiently small end neighborhood $N_n$ of $a_0$
preserves the separation for all $k$.  Its spectral
projection is nonzero; choose a unit vector $\eta_n$ in its
range.  The neighborhood and vector may depend on $n$.

Let $f$ be a continuous function on the target pencil and
$v$ a target label for an evaluation of the composite map.
The commuting-composite assertion gives zero between the
two separated $B_1s_nkA_1^{-1}$ cuts.  Compressing the end
factor to $\eta_n$ yields
\[
 R_{\mathcal U}
  (\id\otimes\langle\,\cdot\,\eta_n,\eta_n\rangle)
   \bigl[f(B_2s_nkA_2^{-1}v)\bigr]
 R_{\mathcal L}=0.
\]
Average over $k$.  In the independent start and end
calculus of copy $2$, Haar averaging sends every
$A_2^{-1}v$ to the same probability $m_S$, irrespective
of its end value or of $\eta_n$.  Thus
\begin{equation}
 \label{eq:gluing-averaged-block}
 R_{\mathcal U}
       \left[\int f(B_2s_nz)\,dm_S(z)\right]
 R_{\mathcal L}=0.
\end{equation}

For any prescribed source point $z_0$, choose split
transports, conjugated in the root-generated source
group, with $(s_n)_*m_S\to\delta_{z_0}$.
The usual rank-one contraction holds away from a proper
repelling set of ordinary measure zero, both for smooth
and for local-field root coordinates.  For every realized
homeomorphism $B_2$, the integral in
Equation~\eqref{eq:gluing-averaged-block} converges to
$f(B_2z_0)$.  Its absolute value is bounded by
$\|f\|_\infty$, so the convergence is strong in the
single-copy spectral calculus.  We obtain
$R_{\mathcal U}f(B_2z_0)R_{\mathcal L}=0$.
Compact exhaustion and the separated-block criterion show
that every $B_2$ evaluation commutes with the whole $B_1$
calculus.  Interchanging start and end proves the
corresponding assertion for $A_1,A_2$.

The fibre identities can first be imposed on countable
dense tests and labels and then extended by continuous
source transport with the bases fixed.  The genuine
source amplification of Section~\ref{sec:transfer}, when
needed, allows exactly this fibrewise transport and
descends all the resulting operator identities.
Within a shared panel, the remaining chamber vertices
already commute with the completion slots.  Together
with Section~\ref{sec:locality} at isolated positions,
this proves the final assertion.
\end{proof}

\subsection{Reflection substitutions and global commutation}

\begin{lemma}[Exact simultaneous reflection]
\label{lem:gluing-reflection}
Let $C,D$ be opposite source chambers.  Reflect them
simultaneously across their apartment panels of types
$i,i^*$, obtaining $C',D'$.  Let $\mathcal F_i$ be the
corresponding simultaneous target reflection at types
$p(i),p(i^*)$.  On the open set of opposite target pairs,
the ordered angular kernel $P_CP_D$ is transported
exactly by $\mathcal F_i$ to $P_{C'}P_{D'}$.
The assertion holds also for the joint two-copy
measurements, with opposition required in each copy.
\end{lemma}

\begin{proof}
The type permutation preserves the Coxeter system, and
opposition on types is respected.  Thus $\mathcal F_i$
is the well-defined involution which reflects an
opposite pair in its common apartment.  On the open
Bruhat orbit it is smooth at real places and continuous
analytic at finite places.  Equivalently, each reflected
chamber is the unique completion of the prescribed panel
which is not opposite to the other original chamber.

We prove the assertion first for smooth compact tests
inside opposition at the real places, and finite clopen
tests with slack in compact opposition patches at the
finite places.  Ordered integrals of the former are
defined by summable separated expansions; the latter
are finite sums.  Insert the two new slots in the order
\[
 P_C\,P_{C'}\,P_{D'}\,P_D.
\]
The first adjacent pair and the last adjacent pair
commute by Lemma~\ref{lem:gluing-two-copy} and the
isolated conclusion of Section~\ref{sec:locality}.
To constrain the $C'$ value, interchange $D,D'$ and
insert the intermediate slots of a gallery from $C'$
to $D$ whose length is one less than the longest
length in the pertinent component.  Every consecutive
pair shares its specified sharp panel, so the
difference of any coordinate of that panel in the two
slots annihilates the ordered kernel with arbitrary
separated multipliers.  The same is then true for the
allowed joint test functions by their convergent
expansions.

At a real place these shared-panel equalities define
a regular submanifold: introduce the chamber variables
successively, each time pulling back a panel equality
along a submersion onto the panel base.  In local
coordinates their differences are independent normal
coordinates.  The integral form of the Taylor formula
in these normal coordinates writes a smooth function
vanishing on the submanifold as a sum of those
differences times smooth functions.  A compact
partition of unity proves annihilation by the entire
smooth vanishing ideal on the testing region.
At a finite place, take finite clopen partitions fine
enough in the compact region.  Shared-panel equalities
then force exact equality of the corresponding
partition values.  Refining the partitions imposes
the same gallery constraint; continuity of the
projection rule on opposition gives its exact
substitution for every clopen test with slack.

On this gallery locus the $C'$ value shares the
designated panel of $C$ and is not opposite to the
$D$ value: the shortened gallery, even with stutters,
cannot have longest Weyl distance.  On original
opposition the panel projection rule therefore
identifies it with the unique reflected completion.
The preceding vanishing-ideal calculation proves
this as a multiplied graph substitution, not just
as a support assertion.  The analogous constraint
on $D'$ follows after initially interchanging
$C,C'$.  The two individual constraints generate
the simultaneous graph on the open region: in
smooth graph coordinates the two new variables are
determined functions of the old pair.  The finite
partition statement gives the same conclusion at
finite places.  It applies separately to both copy
coordinates and therefore to their joint tests.

On opposition of the new values, the identical
argument, interchanging within the adjacent pairs,
gives the inverse graph substitution.  For a test
supported on a compact subset of original opposition,
insert a compact cutoff equal to one on its reflected
image.  Apply the forward substitution in the four
slots and then the inverse substitution to integrate
out the old slots.  The resulting ordered integral
equals the integral against $P_{C'}P_{D'}$ of the
transported test.  All cutoffs stay inside their
respective opposition open sets.  At finite places
choose the cutoffs clopen.  This proves the asserted
identity without crossing the singular boundary of
opposition.  A single isolated rank-one component
already has global commutation and needs no shortened
gallery argument.
\end{proof}

\begin{proposition}[Global joint angular commutation]
\label{prop:gluing-global-commutation}
All chamber measurements, in both normal copies,
commute.  For opposite source chambers the joint
target law is supported on opposition in each copy.
\end{proposition}

\begin{proof}
Apply Lemma~\ref{lem:gluing-reflection} successively
along a reduced longest word, moving the opposite
chamber simultaneously.  In the source apartment
the final pair is $(D,C)$.  The permuted word is
also reduced and longest, so the target pair is
likewise interchanged.  Thus the ordered kernels
$P_CP_D$ and $P_DP_C$, in the same target variables,
have equal integrals for all smooth or clopen compact
tests in opposition.  Approximation by such tests
shows that
\[
 P_C(U)P_D(L)=P_D(L)P_C(U)
 \quad\text{when}\quad
 \overline U\times\overline L
       \subset\operatorname{Opp}(B_T\times B_T).
\]
The same holds for every Borel subcut of these
patches, by spectral regularity.  We first use
this conclusion in one copy.

Let $R_{C,D}$ be the range join of these transverse
products.  Each product $R=P_C(U)P_D(L)$ is a
projection.  It reduces both full PVMs: for an
arbitrary Borel $A$, the subrectangle
$(A\cap U)\times L$ is still transverse, so
$P_C(A)R=RP_C(A)$, and the same argument applies
to $P_D$.  The two restrictions to $\ran R$
commute, again by the subrectangle assertion.
These properties pass to the range join.

The projection $R_{C,D}$ is target invariant,
since opposition is target invariant, and belongs
to the finite algebra $\mathcal N$.  It is fixed
by a source split group fixing $C,D$ and escaping
in every noncompact factor.  The mixing statement
of Section~\ref{sec:setup} applied to
$R_{C,D}-\tau(R_{C,D})1$ makes it zero.  Thus
$R_{C,D}$ is scalar.  It is nonzero: if it vanished
for one opposite source pair, source covariance
and transitivity on the opposite-pair orbit would
make it vanish for every such pair.  Fix nonempty
target patches $U,L$ with compact transverse
product.  For fixed source $C$, almost every
source $D$ is opposite to $C$, so
$P_C(U)P_D(L)=0$ for a conull set of $D$.
The conull-label angular fullness of
Section~\ref{sec:transfer} says that these
$P_D(L)$ ranges span, giving $P_C(U)=0$.
Fullness for $U$ is a contradiction.  Therefore
$R_{C,D}=1$.  In particular the one-copy pair
law, now commuting, gives full mass to opposition.

For clarity, the same scalar-join argument for
the joint copies has a separate nonzeroness step.
Choose a nonzero transverse product
$P_C^1(U)P_D^1(L)$ in copy $1$.  In copy $2$
the preceding paragraph gives a commuting joint
law supported on opposition.  Exhaust this
opposition set by finite disjoint unions of
transverse rectangles, using a countable
rectangle cover and finite refinements.  The
associated sums of products
$P_C^2(U_a)P_D^2(L_a)$ increase strongly to $1$.
Insert these sums between the two copy-$1$
cuts.  Their sandwiched sums converge strongly
to the chosen nonzero product.  Each summand is
\[
 \bigl(P_C^1(U)P_C^2(U_a)\bigr)
 \bigl(P_D^2(L_a)P_D^1(L)\bigr),
\]
because same-chamber copies commute.  This is
exactly a joint transverse product in both
coordinates.  At least one such product is
nonzero.  The joint reflection identity and
scalar-join argument now show that their join
is $1$, so both joint chamber PVMs commute
globally and their law is supported on joint
opposition.  Only separated products have
been inserted; no discontinuous joint cutoff
has been applied to noncommuting slots.

Finally approximate any source pair $(C,D)$
by opposite pairs.  Local source sections and
strong source transport give strong convergence
of each angular operator tested by a continuous
target function.  Uniform boundedness passes
commutation to the limit.  Continuous target
tests generate the spectral algebras, so the
assertion extends to all Borel cuts and to
all source chamber pairs.
\end{proof}

\subsection{From panel bijections to chamber bijections}

\begin{proposition}[Both directions of measure-class preservation]
\label{prop:gluing-chamber-bijections}
In the classical spectral law of one copy, almost every
chamber-map class has a measurable representative
$\phi:B_S\to B_T$ which is a bijection between conull
sets and satisfies
\[
 \phi_*m_S\sim m_T,
\]
where $m_S,m_T$ are the ordinary compact chamber
probabilities.  The same assertion holds for each
normal copy in their joint law.
\end{proposition}

\begin{proof}
Proposition~\ref{prop:gluing-global-commutation}, a
faithful scalar spectral state, and continuity in
measure of the chamber evaluations give jointly
measurable random maps exactly as in the proof of
Proposition~\ref{prop:gluing-homogeneous-pencil}.
On almost every panel, for almost every realization,
the restriction is a nonsingular homeomorphism
with a measurable inverse.  At real non-isolated
positions this is
Proposition~\ref{prop:gluing-homogeneous-pencil};
at finite positions it is Section~\ref{sec:finite};
at isolated positions it is Section~\ref{sec:locality}.
The continuous versions, their inverse maps, and
their evaluations are Borel in the endpoint
charts.  The measure statement follows in all
cases from the respective smooth or analytic
homogeneous identification.  Panel incidence
is preserved with the permutation $p$.

Choose a reduced word $(i_1,\ldots,i_m)$ for the
longest Weyl element of the full product.
Sample $C_0$ with ordinary compact probability,
and successively sample $C_j$ in the
$i_j$-panel of $C_{j-1}$ with its compact
rotation probability.  These transition
probabilities preserve the ordinary chamber
probability: this is disintegration of the
compact-invariant chamber--panel incidence
measure.  Thus each encountered chamber and
panel has its ordinary measure class.
Fubini over the spectral law and the finite
sequence of choices shows that, for almost
every realization and almost every starting
chamber, every needed panel map is good for
almost every path prefix.  Its continuous
version agrees with $\phi$ at the current
chamber and at almost every next choice.
The panel probabilities have no atoms, so
nonstuttering has full measure, and the
injective panel maps preserve nonstuttering.

Fix such a realization and starting chamber
$C_0$.  The resulting map on gallery choices
starts at $\phi(C_0)$ and has a triangular
form: the map on the next panel is determined
by the already mapped prefix.  We verify
that this map preserves measure classes in
both directions.  Suppose this has been
proved through step $j$, with prefix map
$\Phi_j$ and measurable inverse.  Let
$\mu_x$ and $\nu_{\Phi_j(x)}$ be the ordinary
conditional probabilities on the next
source and target pencils.  Outside a
null set of source prefixes there is a
measurable bijection $h_x$ satisfying
\[
 (h_x)_*\mu_x\sim\nu_{\Phi_j(x)}.
\]
By the induction hypothesis, the exceptional
prefix set has null image on the target
side as well.  For a measurable set $A$
of extended target prefixes, conditional
integration gives
\[
 \mu_{j+1}(\Phi_{j+1}^{-1}A)
   =\int \mu_x
       \bigl(h_x^{-1}(A_{\Phi_j(x)})\bigr)
       \,d\mu_j(x).
\]
If $\nu_{j+1}(A)=0$, its conditional sections
are $\nu_{\Phi_j(x)}$-null at almost every
prefix.  Fibre absolute continuity and
$\Phi_{j*}\mu_j\ll\nu_j$ make the displayed
integral zero.  Conversely, if the integral
is zero, nonnegativity makes its integrand
zero almost everywhere.  The reverse fibre
absolute continuity and
$\nu_j\ll\Phi_{j*}\mu_j$ imply
$\nu_{j+1}(A)=0$.  This proves equivalence
of the extended measures.  The inverse
first recovers the prefix through
$\Phi_j^{-1}$ and then uses $h_x^{-1}$;
it is measurable and is inverse almost
everywhere in both directions.  The empty
prefix begins the induction.  We have
therefore proved a measure-class
isomorphism of the full gallery spaces.

The endpoint map identifies the generic
galleries of this fixed reduced longest
type from $C_0$ with its open opposition
cell.  The intermediate chambers are
unique: they are recovered backwards by
successive projections along the descents
of the reduced word.  In root coordinates
this is the ordered unipotent Bruhat
factorization \cite{Borel1991,Tits1964}.
At a real place it is a
smooth coordinate isomorphism with
nonvanishing Jacobian; at a finite place
it is an analytic coordinate isomorphism
in the relative-root coordinates, hence
preserves the local Haar class in both
directions \cite{BruhatTits1972,BruhatTits1984}.
The iterated compact panel
probabilities have exactly these local
classes.  Products of the local
identifications give the assertion for
the full group.  The complementary
Bruhat cells have ordinary measure zero.
The same statements hold for the target
word $(p(i_1),\ldots,p(i_m))$, since $p$
preserves the Coxeter system.

Passing the gallery isomorphism through
these two endpoint identifications gives
a measurable bijection between conull
subsets of $B_S$ and $B_T$ preserving
both null directions.  On almost every
source gallery it agrees at the endpoint
with $\phi$, so it is a version of the
original chamber-map class.  This proves
both essential bijectivity and
$\phi_*m_S\sim m_T$.  Normal copying
preserves all one-copy almost-sure
statements, giving the assertion for
both entries of the joint law.
\end{proof}
